\documentclass{amsart}
% For dealing with references we use the comment environment
\usepackage{verbatim}
\newenvironment{reference}{\comment}{\endcomment}
%\newenvironment{reference}{}{}
\newenvironment{slogan}{\comment}{\endcomment}
\newenvironment{history}{\comment}{\endcomment}

% For commutative diagrams we use Xy-pic
\usepackage[all]{xy}

% We use 2cell for 2-commutative diagrams.
\xyoption{2cell}
\UseAllTwocells

% We use multicol for the list of chapters between chapters
\usepackage{multicol}

% This is generally recommended for better output
\usepackage{lmodern}
\usepackage[T1]{fontenc}

% For cross-file-references

% Package for hypertext links:
\usepackage{hyperref}

% For any local file, say "hello.tex" you want to link to please

% Theorem environments.
%
\theoremstyle{plain}
\newtheorem{theorem}[subsection]{Theorem}
\newtheorem{proposition}[subsection]{Proposition}
\newtheorem{lemma}[subsection]{Lemma}

\theoremstyle{definition}
\newtheorem{definition}[subsection]{Definition}
\newtheorem{example}[subsection]{Example}
\newtheorem{exercise}[subsection]{Exercise}
\newtheorem{situation}[subsection]{Situation}

\theoremstyle{remark}
\newtheorem{remark}[subsection]{Remark}
\newtheorem{remarks}[subsection]{Remarks}

\numberwithin{equation}{subsection}

% Macros
%
\def\lim{\mathop{\mathrm{lim}}\nolimits}
\def\colim{\mathop{\mathrm{colim}}\nolimits}
\def\Spec{\mathop{\mathrm{Spec}}}
\def\Hom{\mathop{\mathrm{Hom}}\nolimits}
\def\Ext{\mathop{\mathrm{Ext}}\nolimits}
\def\SheafHom{\mathop{\mathcal{H}\!\mathit{om}}\nolimits}
\def\SheafExt{\mathop{\mathcal{E}\!\mathit{xt}}\nolimits}
\def\Sch{\mathit{Sch}}
\def\Mor{\mathop{\mathrm{Mor}}\nolimits}
\def\Ob{\mathop{\mathrm{Ob}}\nolimits}
\def\Sh{\mathop{\mathit{Sh}}\nolimits}
\def\NL{\mathop{N\!L}\nolimits}
\def\CH{\mathop{\mathrm{CH}}\nolimits}
\def\proetale{{pro\text{-}\acute{e}tale}}
\def\etale{{\acute{e}tale}}
\def\QCoh{\mathit{QCoh}}
\def\Ker{\mathop{\mathrm{Ker}}}
\def\Im{\mathop{\mathrm{Im}}}
\def\Coker{\mathop{\mathrm{Coker}}}
\def\Coim{\mathop{\mathrm{Coim}}}

% Boxtimes
%
\DeclareMathSymbol{\boxtimes}{\mathbin}{AMSa}{"02}

%
% Macros for moduli stacks/spaces
%
\def\QCohstack{\mathcal{QC}\!\mathit{oh}}
\def\Cohstack{\mathcal{C}\!\mathit{oh}}
\def\Spacesstack{\mathcal{S}\!\mathit{paces}}
\def\Quotfunctor{\mathrm{Quot}}
\def\Hilbfunctor{\mathrm{Hilb}}
\def\Curvesstack{\mathcal{C}\!\mathit{urves}}
\def\Polarizedstack{\mathcal{P}\!\mathit{olarized}}
\def\Complexesstack{\mathcal{C}\!\mathit{omplexes}}
% \Pic is the operator that assigns to X its picard group, usage \Pic(X)
% \Picardstack_{X/B} denotes the Picard stack of X over B
% \Picardfunctor_{X/B} denotes the Picard functor of X over B
\def\Pic{\mathop{\mathrm{Pic}}\nolimits}
\def\Picardstack{\mathcal{P}\!\mathit{ic}}
\def\Picardfunctor{\mathrm{Pic}}
\def\Deformationcategory{\mathcal{D}\!\mathit{ef}}

\setlength{\emergencystretch}{3em}
\begin{document}
\title[Stacks Project, Tag 0D83]{320 --- Existence and isomorphism of Postnikov systems under Hom vanishing}
\maketitle
\noindent Copyright (C) 2005--2025 Johan de Jong. Stacks Project authors. Source: \href{https://stacks.math.columbia.edu/tag/0D83}{Tag 0D83}.

\noindent Editorial difficulty note (not author text). Difficulty H2: The proof inductively extends a totalization tower and identifies its obstruction as a composite into an earlier cone. Iterated distinguished triangles and the shifted Hom bounds annihilate that obstruction; the stronger bounds then lift the identity to an isomorphism of systems. This is an obstruction-theoretic existence theorem for totalizing complexes in abstract triangulated categories..

\noindent Editorial provenance note (not author text). History: excerpt from revision \texttt{a0444\allowbreak 6e57e\allowbreak c1fbc\allowbreak 252a8\allowbreak 71afc\allowbreak ec775\allowbreak 2fb28\allowbreak 07b14}, derived.tex, lines 12319--12370. Reference presentation was adapted; original-author context and core support blocks were retained or added. The mathematical wording is unchanged. Licensed under GNU FDL 1.2 or later; no invariant sections or cover texts. See ../licenses/COPYING. Context blocks reproduce author definitions and supporting results; they are not additional counted problems.

\begin{definition}
\label{definition-postnikov-system}
Let $\mathcal{D}$ be a triangulated category. Let
$$
X_n \to X_{n - 1} \to \ldots \to X_0
$$
be a complex in $\mathcal{D}$. A {\it Postnikov system} is defined
inductively as follows.
\begin{enumerate}
\item If $n = 0$, then it is an isomorphism $Y_0 \to X_0$.
\item If $n = 1$, then it is a choice of an isomorphism $Y_0 \to X_0$ and
a choice of a distinguished triangle
$$
Y_1 \to X_1 \to Y_0 \to Y_1[1]
$$
where $X_1 \to Y_0$ composed with $Y_0 \to X_0$ is the given morphism
$X_1 \to X_0$.
\item If $n > 1$, then it is a choice of a Postnikov system
for $X_{n - 1} \to \ldots \to X_0$ and a choice of a distinguished
triangle
$$
Y_n \to X_n \to Y_{n - 1} \to Y_n[1]
$$
where the morphism $X_n \to Y_{n - 1}$ composed with
$Y_{n - 1} \to X_{n - 1}$ is the given morphism $X_n \to X_{n - 1}$.
\end{enumerate}
Given a morphism
\begin{equation}
\label{equation-map-complexes}
\vcenter{
\xymatrix{
X_n \ar[r] \ar[d] &
X_{n - 1} \ar[r] \ar[d] &
\ldots \ar[r] &
X_0 \ar[d] \\
X'_n \ar[r] &
X'_{n - 1} \ar[r] &
\ldots \ar[r] &
X'_0
}
}
\end{equation}
between complexes of the same length in $\mathcal{D}$
there is an obvious notion of a {\it morphism of Postnikov systems}.
\end{definition}

\begin{lemma}
\label{lemma-existence-postnikov-system}
Let $\mathcal{D}$ be a triangulated category.
Let $X_n \to X_{n - 1} \to \ldots \to X_0$ be
a complex in $\mathcal{D}$. If
$$
\Hom(X_i[i - j - 2], X_j) = 0 \text{ for }i > j + 2
$$
then there exists a Postnikov system. If we have
$$
\Hom(X_i[i - j - 1], X_j) = 0 \text{ for }i > j + 1
$$
then any two Postnikov systems are isomorphic.
\end{lemma}

\begin{proof}
We argue by induction on $n$. The cases $n = 0, 1, 2$
follow from Lemma \href{https://stacks.math.columbia.edu/tag/0D81}{lemma postnikov system small cases (Tag 0D81)}.
Assume $n > 2$.
Suppose given a Postnikov system for the complex
$X_{n - 1} \to X_{n - 2} \to \ldots \to X_0$.
The only obstruction to extending this to a Postnikov system
of length $n$ is that we have to find a morphism
$X_n \to Y_{n - 1}$ such that the composition
$X_n \to Y_{n - 1} \to X_{n - 1}$ is equal to the given map
$X_n \to X_{n - 1}$. Considering the distinguished triangle
$$
Y_{n - 1} \to X_{n - 1} \to Y_{n - 2} \to Y_{n - 1}[1]
$$
and the associated long exact sequence coming from this
and the functor $\Hom(X_n, -)$
(see Lemma \href{https://stacks.math.columbia.edu/tag/0149}{lemma representable homological (Tag 0149)})
we find that it suffices to show that the composition
$X_n \to X_{n - 1} \to Y_{n - 2}$ is zero.
Since we know that $X_n \to X_{n - 1} \to X_{n - 2}$ is zero
we can apply the distinguished triangle
$$
Y_{n - 2} \to X_{n - 2} \to Y_{n - 3} \to Y_{n - 2}[1]
$$
to see that it suffices if $\Hom(X_n, Y_{n - 3}[-1]) = 0$.
Arguing exactly as in the proof of
Lemma \href{https://stacks.math.columbia.edu/tag/0D82}{lemma maps postnikov systems vanishing (Tag 0D82)} part (1)
the reader easily sees this follows from the condition
stated in the lemma.

\medskip\noindent
The statement on isomorphisms follows from the existence of a map
between the Postnikov systems extending the identity on the complex
proven in Lemma \href{https://stacks.math.columbia.edu/tag/0D82}{lemma maps postnikov systems vanishing (Tag 0D82)} part (2)
and Lemma \href{https://stacks.math.columbia.edu/tag/014A}{lemma third isomorphism triangle (Tag 014A)} to show all the maps are
isomorphisms.
\end{proof}
\end{document}
